\chapter{Against Kyubey}
Behind every story is at least one storyteller. As I argued back in Chapter 5, there is reason to equate Kyubey at least partially with the creators of Madoka Magica, and in particular with Urobuchi himself. It is, therefore, interesting that despite Homura's protestations that she has become Evil at the end of the film, in much the same way Madoka became Hope at the end of the series, two earlier scenes imply that we should look past this to find the true villain, and thereby the true tyrant being rebelled against.

Relatively early in the film, after Homura has begun to realize that the magical girls' memories are being tampered with, she discovers that the space within which they are sealed resembles a witch's labyrinth. Recognizing Bebe as the witch Charlotte, she concludes it is her doing. There follows a scene in which Madoka, Mami, and Homura sit around Mami's table, which (as it was in the series) is shaped like an extended isosceles triangle. Madoka makes a comment which Homura then uses as an opening to begin questioning Mami about where Bebe comes from, to which Bebe responds with confusion and concern. Throughout the scene, Homura is attempting to confirm her suspicions that Bebe is the creator of the sealed space, and therefore the villain of this story-but the table is pointing directly at Kyubey throughout, as shown both in high shots that depict the entire group from above, and even more tellingly in shots along the center-line of the table, which emphasize it pointing past Bebe to Kyubey.

Correspondingly, near the end of the film, the magical girls initially appear to be fighting Homura's army of familiars. However, once they succeed in shattering the barrier, we are able to see massive Kyubeys looking down through their isolation field, the true prison containing the magical girls (Homura included)-once again, looking past Homura to see the true enemy.

In the previous chapter, I briefly discussed E.T.A. Hoffman's "The Nutcracker and the Mouse Prince." Here it can be discussed more fully, because the final battle in many ways mimics elements of the battle between Marie's toys and the Mouse King early in that story: the toys are commanded by a broken-jawed creature, march forth from a cabinet, and battle an evil mouselike creature. Given his seven heads with seven crowns (features explicitly mentioned in the story), the Mouse King distinctly recalls the seven-headed, ten-crowned beast of Revelation 13, but in terms of Madoka Magica the many heads in one body serve as a reflection of Kyubey's single mind in many bodies. Just as the Mouse King is the enemy of Marie (Mary) and the Beast is an enemy of God, Kyubey is an enemy of Madoka (Kanon), and Homura acting as her knight is therefore the Nutcracker himself. Finally, the Mouse King demands to feed on Marie's candies and toys as a ransom for sparing the Nutcracker, much as Kyubey demands the magical girls sacrifice their lives and souls to feed his need for energy-food being the most basic source of energy for a living animal. (Including, notably, a sort of edible doll that is Marie's favorite and described quite similarly to the Nutcracker's true form/the "real" Drosselmeier's nephew, particularly with the phrase "rosy-cheeked." If the Nutcracker is Homura as she is now, the rosy-cheeked sugar doll is the original Homura, the weak girl that Madoka rescued only to unintentionally draw her into Kyubey's orbit.)

Most interestingly in this light, in "The Nutcracker and the Mouse King" the figure of the storyteller is deeply ambiguous. Drosselmeier, who gives the Nutcracker to Marie as a gift and tells her the story of his origins and curse, is repeatedly implied to be untrustworthy or allied with the Mouse King. Most obviously, he either permits or causes the clock to strike, summoning the Mouse King to invade Marie's room and forcing the Nutcracker to raise the army of toys. Later, although Drosselmeier seems rather kindly in telling the injured Marie his story to cheer her up, he dismisses as dreams and fantasies her claims to have witnessed the battling toys or the Nutcracker's final slaying of the Mouse King. 

Drosselmeier also inserts himself as a character into the story of the Nutcracker's curse, much as Kyubey inserts himself into Homura's witch barrier (which is only possible because it exists within Kyubey's isolation field). As he tells it, he is more or less the hero of the story-within-a-story, going on a quest to save the Princess Pirlipat from a curse. However, the Drosselmeier of the inner story is equally readable, especially once the reader begins to suspect the "real" Drosselmeier, as a villain, a classic fairy-tale wicked uncle who gets himself into trouble, and sacrifices his own nephew to save himself. We have only Drosselmeier's word for it that the ritual went wrong-perhaps it was always intended to merely transfer the curse from the Princess onto the Nutcracker? His story cannot be trusted-while Marie's experiences are always narrated in the same way as the straightforwardly "real" events of the story, despite occurring in times and places that make it possible to dismiss them as dreams, Drosselmeier's story is presented by a character, as a fictional tale to entertain a child. Within the context of the "The Nutcracker and the Mouse King," Drosselmeier's tale of the Nutcracker is less likely to be "true" than Marie's dreams. The ending further implies her dreams are "realer" than his stories, as Marie is able to leave for the Nutcracker's world, and he describes Drosselmeier with contempt.

Whether Drosselmeier is more villainous than he presents himself is up to the reader's preference and interpretation. Certainly, however, there must be some reason that the Nutcracker, when Marie enters his realm and cures him (just as Madoka reaches through the cabinet and restores Homura from being a witch), is openly dismissive of Drosselmeier, warning Marie that he cannot deliver on the clockwork lake he promised to make for her, and that her power of creation is greater than his. And in the end, the happy ending for Marie is for her and the Nutcracker to escape "reality" into the Nutcracker's realm.

While the details of the plot do not mesh, in particular the order of events, the character parallels are quite clear. Homura is responsible for creating the imagery within her barrier, and the recurring theme of the Nutcracker in association with her suggests that she sees herself as that figure, taking the Princess' curse onto himself and living for the possibility of one day being rescued by Marie. Madoka is both the Princess and Marie (and it is notable that, in the Nutcracker's realm, Marie sees her reflection in a pool as the Princess-the two are quite possibly the same character in the Hoffmann story as well), while Kyubey combines elements of the Mouse King and Drosselmeier.

These mappings do not necessarily reflect the true natures of the characters, it should be noted, as they are drawn from the imagery of the witch's barrier, which has before been established as a hallucinatory maze that reflects the witch's state of mind, imposed on and replacing reality. But it gives a great deal of insight into Homura's motivations; the happy ending for Marie is to escape reality, not to some nirvana or heaven, but to the Nutcracker's world, where they can live together. Homura creates her world to be with Madoka, her Marie, in a dark and desperate attempt to complete her version of the Nutcracker. She is, more than ever, that dark and ugly being-and the only cure is for someone to love her as she is. For, in other words, Madoka to learn what Homura has done and has become, and still not give up on her.

But Drosselmeier may well have the last laugh. As a storyteller within a story, he is necessarily therefore at least in part a reflection of his own storyteller, Hoffmann-and Marie and the Nutcracker are still within Hoffmann's domain at the end of the story, as witness that he is able to keep writing about them until he chooses to stop. Kyubey is Drosselmeier is Hoffman is Urobuchi; he appears to have been defeated, but he seemed defeated at the end of the series as well. The final shot of his staring eye reminds us: as long as we are watching, he is watching. As long as we are willing to accept further stories of Madoka, Homura, and all the rest, Kyubey retains the power to threaten them.
